\section{\texorpdfstring{Derived Categories:\allowbreak{} essentially constant systems,\allowbreak{} original 12380–\allowbreak{}12604}{Derived  essentially constant  original }}\label{proofs-12380-12604-derived-categories-essentially-constant-systems-original-1238012604}

Authority:\allowbreak{} Stacks Project commit a04446e5\allowbreak{}7ec1fbc2\allowbreak{}52a871af\allowbreak{}cec7752f\allowbreak{}b2807b14, derived.tex,\allowbreak{} SHA-256 EFDA4CB3\allowbreak{}61DD4090\allowbreak{}9D199116\allowbreak{}1EAE716E\allowbreak{}58B6CF9E\allowbreak{}79380971\allowbreak{}502E1DA1\allowbreak{}2396D402. This is a complete editorial proof supplement for the terminal mathematical section. DERIVED-NEW-044 repairs a finite-stage proof gap without changing the theorem. DERIVED-UNDER-033 records the finite cohomological-degree bound proved below; no novelty is claimed.

\subsection{\texorpdfstring{1. Essential constancy,\allowbreak{} its splitting,\allowbreak{} and the actual pro-category maps}{1. Essential  its  and the actual pro-category maps}}\label{proofs-12380-12604-essential-constancy-its-splitting-and-the-actual-pro-category-maps}

For a directed inverse system write \(t_{ji}:A_j\to A_i\) for \(j\geq i\). The source definition of essential constancy gives an object \(A\),\allowbreak{} a cone \(q_i:A\to A_i\),\allowbreak{} and \(p:A_{i_0}\to A\) with \(pq_{i_0}=1_A\),\allowbreak{} such that for each \(i\) some \(j\geq i,i_0\) satisfies \[
t_{ji}=q_i p t_{j i_0}. \tag{1.1}
\] This is the inverse-system specialization of Categories 2945–\allowbreak{}2955. It is stronger than mere existence of a limit.

For \(i\geq i_0\),\allowbreak{} put \(p_i=p t_{i i_0}\); then \(p_iq_i=1_A\). Complete \(p_i\) to a triangle \[
Z_i\xrightarrow{z_i}A_i\xrightarrow{p_i}A\xrightarrow{\epsilon_i}Z_i[1].
\] Consecutive composites vanish,\allowbreak{} so \(\epsilon_i=\epsilon_ip_iq_i=0\). The split-triangle proof (original Derived 699–\allowbreak{}726)\allowbreak{} makes \(z_i\oplus q_i:Z_i\oplus A\to A_i\) an isomorphism. For any \(T\),\allowbreak{} exact Hom gives \(0\to\operatorname{Hom}(T,Z_i)\to\operatorname{Hom}(T,A_i)
 \to\operatorname{Hom}(T,A)\to0\); the initial zero follows from \(\epsilon_i[-1]=0\). Thus \(z_i\) is an actual categorical kernel of \(p_i\). The identities \(p_i t_{ji}=p_j\) and \(t_{ji}q_j=q_i\) show the transition matrix is \(\operatorname{diag}(1_A,z_{ji})\) in the ordered decomposition \(A\oplus Z_i\). Equation (1.1)\allowbreak{} says \(z_{ji}=0\) for a sufficiently late \(j\) at each fixed \(i\). This proves the source\textquotesingle s forward implication without assuming arbitrary idempotents split in the triangulated category:\allowbreak{} the exhibited retraction has a complement by its triangle.

Conversely,\allowbreak{} suppose the displayed decompositions and transition matrices of the source lemma exist on a tail. The inclusions of \(A\) give a cone there,\allowbreak{} and hence on the whole system by composing to each earlier index; independence follows by taking a common later index. The projection at the first index has composite identity on \(A\). For each \(i\),\allowbreak{} choose \(j\) killing \(Z_j\to Z_i\); then the transition factors through \(A\) exactly as in (1.1)\allowbreak{}. This is essential constancy. For a compatible cone from any \(T\),\allowbreak{} its component into \(Z_i\) is zero,\allowbreak{} because it factors through a transition \(Z_j\to Z_i\) which is zero. Its components into \(A\) are equal. Consequently the actual limit is \(A\),\allowbreak{} with the specified cone maps.

For clarity about the subsequent pro-category argument,\allowbreak{} let \[
F_X(T)=\operatorname{colim}_i\operatorname{Hom}(X_i,T).
\] A pro-map \(a:X\to Y\) consists of compatible classes \([a_j]\in\operatorname{colim}_i\operatorname{Hom}(X_i,Y_j)\). It induces \(F_Y\to F_X\):\allowbreak{} the class of \(u:Y_j\to T\) is sent to the class of \(u a_j\). Refining indices proves this is independent of both representatives; compatibility in \(j\) proves it respects the colimit. It is natural in \(T\) and respects composition. Conversely,\allowbreak{} evaluate a natural transformation \(F_Y\to F_X\) at the class of \(1_{Y_j}\); the resulting classes give a pro-map,\allowbreak{} and naturality recovers every class \(u\). These constructions are inverse. Equivalently,\allowbreak{} \[
\operatorname{Nat}(F_Y,F_X)
 =\lim_j\operatorname{colim}_i\operatorname{Hom}(X_i,Y_j)
 =\operatorname{Hom}_{\operatorname{Pro}(\mathcal D)}(X,Y). \tag{1.2}
\] This proves the exact contravariant fully faithful comparison cited by the source,\allowbreak{} rather than using only the name of the copresheaf construction.

An isomorphism from \(X\) to a constant pro-object \(A\) gives an actual cone \(A\to X_i\) and a map \(X_{i_0}\to A\) representing its inverse. The constant composite is literally \(1_A\); equality of the other composite with the pro-identity is precisely eventual equality (1.1)\allowbreak{} at each index. Thus pro-isomorphism to a constant system and the source\textquotesingle s essential constancy are equivalent.

\subsection{\texorpdfstring{2. DERIVED-NEW-044:\allowbreak{} the needed additional tail in the triangle comparison}{2.  the needed additional tail in the triangle comparison}}\label{proofs-12380-12604-derived-new-044-the-needed-additional-tail-in-the-triangle-comparison}

Suppose \[
A_n\xrightarrow{a_n}B_n\xrightarrow{b_n}C_n
 \xrightarrow{d_n}A_n[1] \tag{2.1}
\] is an inverse sequence of triangles and the outer systems are essentially constant. After taking a common tail,\allowbreak{} Section 1 gives decompositions \(A_n=A\oplus A'_n\),\allowbreak{} \(C_n=C\oplus C'_n\),\allowbreak{} with diagonal transitions and pro-zero primed systems. Write the actual connecting map in blocks:\allowbreak{} \[
d_n=
\begin{pmatrix}
\delta_n&\beta_n\\
\gamma_n&\eta_n
\end{pmatrix}
:C\oplus C'_n\longrightarrow A[1]\oplus A'_n[1]. \tag{2.2}
\] For \(m\geq n\),\allowbreak{} compatibility gives \[
\delta_m=\delta_n,\qquad
\gamma_n=a'_{mn}[1]\gamma_m,\qquad
\beta_m=\beta_n c'_{mn},\qquad
a'_{mn}[1]\eta_m=\eta_n c'_{mn}. \tag{2.3}
\] At each fixed \(n\),\allowbreak{} choose \(m\) killing \(a'_{mn}\); then \(\gamma_n=0\). Thus the assertion in original 12450–\allowbreak{}12451 that \(C\) maps into \(A[1]\) by a fixed map \(\delta\) is correct.

However,\allowbreak{} the projections \(C_n\to C\) and \(A_n[1]\to A[1]\) commute with the connecting maps only when \(\beta_n=0\). This does not follow at the first index merely from \(\gamma_n=0\). Fix the first remaining index \(n_0\),\allowbreak{} and choose \(N\geq n_0\) with \(c'_{N n_0}=0\). Equation (2.3)\allowbreak{} gives \(\beta_N=0\),\allowbreak{} and then \(\beta_n=0\) for every \(n\geq N\). Passing to this additional tail makes the needed square commute. This is the missing finite-stage justification before the TR3 step at original 12456–\allowbreak{}12461.

The failure at an unadjusted first index is explicit. In \(D^b(k)\),\allowbreak{} put \(K=k[0]\). Let \(A_n=K\) with identity transitions; let \(C_1=K[1]\) and \(C_n=0\) for \(n\geq2\); let \(B_1=0\) and \(B_n=K\) for \(n\geq2\). At the first index use \[
K\longrightarrow0\longrightarrow K[1]\xrightarrow{1}K[1],
\] and at every later index use \(K\xrightarrow{1}K\to0\to K[1]\). The transition from index two to one has components \(1_K,0,0\); every square commutes,\allowbreak{} including the shifted last square. The outer systems have values \(K\) and \(0\). Their constant triangle is \(K\xrightarrow{1}K\to0\to K[1]\). A comparison from the first triangle to this one inducing the identity on \(A_1=K\) would force \(0=1_K\) in the first square. There is no such map. At index two it does exist. Thus the proof needs the additional tail; the theorem,\allowbreak{} which asks for some index,\allowbreak{} is unchanged.

After the adjustment,\allowbreak{} choose a distinguished triangle \[
A\xrightarrow{a}B\xrightarrow{b}C\xrightarrow{\delta}A[1].
\] Apply TR3 to the square on \(C_N,A_N[1]\),\allowbreak{} then rotate back,\allowbreak{} obtaining a triangle map from (2.1)\allowbreak{} at \(N\) to this fixed triangle. For \(n\geq N\),\allowbreak{} compose with the given transition triangle maps to obtain compatible maps \(p^B_n:B_n\to B\).

For every \(T\),\allowbreak{} the induced long exact Hom sequences give a natural diagram \[
\operatorname{Hom}(C[s],T)\to\operatorname{Hom}(B[s],T)
 \to\operatorname{Hom}(A[s],T)\to\operatorname{Hom}(C[s-1],T)
\] over the corresponding filtered colimits of the groups from (2.1)\allowbreak{},\allowbreak{} for every integer \(s\). The colimit sequence is exact:\allowbreak{} for an element whose image is zero,\allowbreak{} represent it at one stage,\allowbreak{} advance to a stage where its image is literally zero,\allowbreak{} and use exactness at that stage. The outer comparisons are isomorphisms because the primed systems are pro-zero. Applying the five lemma to the five consecutive terms surrounding the middle term makes the middle comparison an isomorphism. These are maps of functors in \(T\),\allowbreak{} not just group isomorphisms selected separately. Section 1 identifies the resulting comparison as a pro-isomorphism \(B_\bullet\to B\),\allowbreak{} hence as essential constancy with value \(B\).

For completeness,\allowbreak{} a pro-isomorphism from a system \(B_n\) to a constant \(B\),\allowbreak{} represented on a tail by the compatible maps \(p^B_n\),\allowbreak{} induces the required isomorphism on its actual limit. Let \(q_n:B\to B_n\) be the inverse cone. Then \(p^B_nq_n=1_B\) at every index where \(p^B_n\) represents the pro-map,\allowbreak{} and \(q_n p^B_m=t_{mn}\) after a sufficiently late \(m\),\allowbreak{} for each fixed \(n\). These are the identities from essential constancy and show its limit cone is \(q_n\). The same reasoning applies to \(A,C\). Thus the final limit clause of the source lemma is proved as well.

The minimal correction adds the calculation killing the components \(C'_n\to A[1]\) and one more increase of the starting index. The original statement,\allowbreak{} its triangle,\allowbreak{} and all receiving hypotheses are preserved.

\subsection{3. Essentially constant bounded cohomology}\label{proofs-12380-12604-essentially-constant-bounded-cohomology}

Suppose every \(E_n\) has cohomology in the same original interval \([a,b]\),\allowbreak{} and every system \(H^i(E_n)\) is essentially constant. If \(a=b\),\allowbreak{} the natural truncation comparisons identify \(E_n\) with \(H^a(E_n)[-a]\),\allowbreak{} compatibly with its transition maps. Thus the system is essentially constant.

For \(a<b\),\allowbreak{} use the actual functorial triangles \[
\tau_{\leq a}E_n\longrightarrow E_n\longrightarrow
\tau_{\geq a+1}E_n\longrightarrow(\tau_{\leq a}E_n)[1].
\] The first system is naturally isomorphic to \(H^a(E_n)[-a]\). The third has cohomology in \([a+1,b]\) and meets the induction hypothesis. Section 2 proves essential constancy of the middle system. This induction ends after the finite interval; no exactness of products,\allowbreak{} enough injectives,\allowbreak{} or existence of an arbitrary derived inverse limit is used.

Let its value be \(E\). Applying any functor \(G\) to the explicit identities (1.1)\allowbreak{} and \(pq=1\) shows that \(G(E_n)\) is essentially constant with value \(G(E)\). In particular,\allowbreak{} \(H^i(E)\) is the original value of \(H^i(E_n)\). This is the omitted cohomology assertion,\allowbreak{} and it does not assume cohomology commutes with arbitrary inverse limits.

\subsection{4. Pro-isomorphisms and the actual cone sequence}\label{proofs-12380-12604-pro-isomorphisms-and-the-actual-cone-sequence}

If \(A_n\to B_n\to C_n\to A_n[1]\) is an inverse sequence of triangles with \(C_n\) pro-zero,\allowbreak{} then both \(\operatorname{colim}_n\operatorname{Hom}(C_n,T)\) and \(\operatorname{colim}_n\operatorname{Hom}(C_n[-1],T)\) are zero. Indeed,\allowbreak{} a representative map out of \(C_n\) becomes zero after precomposition with a zero transition. Exactness gives a natural isomorphism \[
\operatorname{colim}_n\operatorname{Hom}(B_n,T)
 \longrightarrow\operatorname{colim}_n\operatorname{Hom}(A_n,T).
\] The explicit comparison (1.2)\allowbreak{} proves that the original maps \(A_n\to B_n\) give a pro-isomorphism.

Now suppose a strict sequence of maps \(f_n:A_n\to B_n\) has both systems in the same interval \([a,b]\),\allowbreak{} and every \(H^p(f_n)\) is a pro-isomorphism. Choose a cone triangle for each \(f_n\). For each adjacent transition,\allowbreak{} TR3 supplies a third component; define transitions over larger gaps by composition. Because the index is the natural numbers,\allowbreak{} this constructs a coherent inverse sequence of triangles. It does not claim an arbitrary diagram of cones has automatically coherent choices.

The long exact cohomology sequence shows \[
H^p(C_n)=0\quad\hbox{if }p<a-1\hbox{ or }p>b. \tag{4.1}
\] The lower endpoint is \(a-1\),\allowbreak{} not \(a\). For each \(p\),\allowbreak{} put \[
Q_n^p=\operatorname{coker}H^p(f_n),\qquad
J_n^{p+1}=\ker H^{p+1}(f_n).
\] There is a compatible short exact sequence \[
0\longrightarrow Q_n^p\longrightarrow H^p(C_n)
 \longrightarrow J_n^{p+1}\longrightarrow0. \tag{4.2}
\]

Here is the exact reason these two outside systems are pro-zero. For a strict map \(u:X_\bullet\to Y_\bullet\) with pro-inverse,\allowbreak{} the equality \(u u^{-1}=1_Y\) implies that,\allowbreak{} at every target index \(n\),\allowbreak{} some transition \(Y_m\to Y_n\) factors through \(u_n\). It therefore kills the cokernel system at \(n\). Similarly \(u^{-1}u=1_X\) implies that after advancing the source index some transition \(X_k\to X_n\) has the form \(v u_k\) for an actual map \(v:Y_k\to X_n\),\allowbreak{} possibly after refining the representative of the inverse. It kills \(\ker u_k\). These assertions follow directly from eventual equality in the filtered colimit defining each component of a pro-map; they do not use an unproved abelian structure on a pro-category.

Fix \(n\). Choose \(m\geq n\) killing \(Q_m^p\to Q_n^p\),\allowbreak{} and then \(k\geq m\) killing \(J_k^{p+1}\to J_m^{p+1}\). By (4.2)\allowbreak{},\allowbreak{} the image of \(H^p(C_k)\to H^p(C_m)\) lies in \(Q_m^p\); its image at \(n\) is therefore zero. Thus \(H^p(C_n)\) is pro-zero. Together with (4.1)\allowbreak{},\allowbreak{} Section 3 makes \(C_n\) essentially constant with an object \(C\) whose cohomology is zero in every degree. Such an object is zero in the derived category,\allowbreak{} by its definition as the localization at quasi-isomorphisms. Hence \(C_n\) is pro-zero. The preceding paragraph proves the required pro-isomorphism \(A_\bullet\to B_\bullet\).

The converse also holds for these maps:\allowbreak{} a pro-isomorphism remains one after the functor \(H^p\),\allowbreak{} because functors preserve the component maps,\allowbreak{} refinements,\allowbreak{} and inverse equations. Thus,\allowbreak{} in the common bounded interval,\allowbreak{} cohomological pro-isomorphism is an exact detection criterion.

\subsection{\texorpdfstring{5. DERIVED-UNDER-033:\allowbreak{} a finite-degree annihilation bound}{5.  a finite-degree annihilation bound}}\label{proofs-12380-12604-derived-under-033-a-finite-degree-annihilation-bound}

A morphism \(f:E\to F\) in \(D(\mathcal A)\) will be called cohomologically zero when \(H^q(f)=0\) for every integer \(q\). Retain the original degree labels and canonical truncation morphisms.

For each \(q\),\allowbreak{} the map \(\tau_{\leq q}f\) factors through \(\tau_{\leq q-1}F\). To prove this,\allowbreak{} compose it with \(\tau_{\leq q}F\to H^q(F)[-q]\). The restriction of this composite to \(\tau_{\leq q-1}E\) is zero by \(\operatorname{Hom}(D^{\leq q-1},D^{\geq q})=0\). Exact Hom for the truncation triangle of \(E\) makes the composite factor through \(H^q(E)[-q]\). Its induced map on degree-\(q\) cohomology is \(H^q(f)=0\),\allowbreak{} so that intervening morphism between shifted heart objects is zero. Thus the whole composite is zero. Exact Hom for the truncation triangle of \(F\) now gives the asserted factorization. This constructs the actual morphism needed to pass down one filtration step; it is not assumed to be unique.

Let \(S=\{s_1<\cdots<s_r\}\) be a finite set of integers and suppose every object in a chain \[
E_0\xrightarrow{f_1}E_1\xrightarrow{f_2}\cdots
 \xrightarrow{f_r}E_r
\] has cohomology zero outside these exact degrees. If every \(f_j\) is cohomologically zero,\allowbreak{} their composite is zero. Indeed \(E_0\cong\tau_{\leq s_r}E_0\). The factorization just proved makes the first map land in \(\tau_{\leq s_r-1}E_1\). Because every cohomology group in the gap \((s_{r-1},s_r)\) is zero,\allowbreak{} the canonical map \(\tau_{\leq s_{r-1}}E_1\to\tau_{\leq s_r-1}E_1\) is a quasi-isomorphism; use its inverse in the derived category to factor through the former. Naturality of truncation allows composition with \(\tau_{\leq s_{r-1}}f_2\),\allowbreak{} which then factors through \(\tau_{\leq s_{r-1}-1}E_2\). Repeating gives a factorization of the complete composite through \(\tau_{\leq s_1-1}E_r=0\). For \(r=1\) this uses the first factorization directly; if \(S\) is empty,\allowbreak{} every object is already zero. Every skipped degree is retained as an explicitly vanishing cohomology group,\allowbreak{} not removed by a coordinate change.

In particular,\allowbreak{} for a common interval \([a,b]\),\allowbreak{} any composite of \(b-a+1\) cohomologically zero maps is zero. A cohomologically zero endomorphism has \(f^{b-a+1}=0\). More precisely,\allowbreak{} the number of possibly nonzero original cohomological degrees gives the bound \(r\).

For an inverse system supported in \(S\) whose cohomology systems are pro-zero,\allowbreak{} fix an index \(n_0\). At each step choose \(n_{j+1}\geq n_j\) so that the transition kills all \(H^{s_\ell}\),\allowbreak{} \(1\leq\ell\leq r\),\allowbreak{} simultaneously. Such an index exists because finitely many indices have a common upper bound. All other cohomology groups vanish from the outset. The transition \[
E_{n_r}\longrightarrow E_{n_{r-1}}\longrightarrow\cdots
 \longrightarrow E_{n_0}
\] is a composite of \(r\) cohomologically zero maps and hence is zero. This supplies a direct finite procedure proving pro-zero,\allowbreak{} not just an existence statement obtained from an abstract constant value.

Applied to the cones in Section 4,\allowbreak{} there are at most \(b-a+2\) possible degrees,\allowbreak{} namely \([a-1,b]\). At each target index the two-step construction from (4.2)\allowbreak{} can be performed for all these finitely many degrees simultaneously; repeating this for \(b-a+2\) transitions makes the cone transition itself zero. This quantitatively propagates the consequence into the source\textquotesingle s pro-isomorphism criterion.

A finite common range cannot be dropped without a replacement hypothesis. In \(D(k)\),\allowbreak{} let \(E_n\) be the actual complex with \(k\) in every degree \(q\geq n\),\allowbreak{} zero in smaller degrees,\allowbreak{} and all differentials zero. Use the literal inclusions \(E_{n+1}\to E_n\). For every fixed \(q\),\allowbreak{} the system \(H^q(E_n)\) is eventually zero. Yet every transition \(E_m\to E_n\),\allowbreak{} \(m\geq n\),\allowbreak{} is nonzero:\allowbreak{} its cohomology map in degree \(m\) is \(1_k\). Therefore the system is not pro-zero. If it were essentially constant with value \(E\),\allowbreak{} Section 3\textquotesingle s functorial argument would give \(H^q(E)=0\) for all \(q\),\allowbreak{} hence \(E=0\),\allowbreak{} a contradiction. This exhibits the exact limitation of the source result while the finite-degree result proves what does work.

These bounds are editorial consequences of truncation orthogonality,\allowbreak{} with no novelty claim and no hypothesis removed from the source body.

\subsection{6. Remaining physical reports}\label{proofs-12380-12604-remaining-physical-reports}

P02-E0100 (producer line 103)\allowbreak{} and French DERIVED-091 (968)\allowbreak{} both concern original 12551–\allowbreak{}12555. Insert ``Let'' before the displayed inverse system to supply the missing introductory verb. This is exactly existing MC-STK-ERR-0902.

P02-E0101 (104)\allowbreak{} and French DERIVED-092 (969)\allowbreak{} both concern original 12559–\allowbreak{}12560. The head noun ``system'' is singular,\allowbreak{} so ``define'' becomes ``defines''. The p02 report inaccurately names its subject ``This''; the original sentence and French report identify the actual subject. The replacement itself is correct and is exactly existing MC-STK-ERR-0903. Both units are replayed against authority and current bytes; no duplicate operation is added.

These four physical reports form groups DERIVED-RECON-102 and DERIVED-RECON-103. All 176 chapter reports have now been individually accounted for:\allowbreak{} 160 already corrected,\allowbreak{} thirteen missing copyedit/\allowbreak{}notation reports,\allowbreak{} two missing mathematical-index reports,\allowbreak{} and one metadata-only resolution. The independent proof gap in Section 2 is not counted as a received report.

\subsection{\texorpdfstring{7. Receiving uses,\allowbreak{} completed scope,\allowbreak{} and next work}{7. Receiving  completed  and next work}}\label{proofs-12380-12604-receiving-uses-completed-scope-and-next-work}

The literal-reference search of current root chapter sources routes the two-out-of-three theorem to duality.tex 8222 and 8360,\allowbreak{} and the cohomology/\allowbreak{}pro-isomorphism results to more-algebra.tex 27867,\allowbreak{} 27872,\allowbreak{} 29719,\allowbreak{} 30341,\allowbreak{} 30346 and 30389.

In duality.tex 8218–\allowbreak{}8224 the inputs are inverse sequences of restricted triangles with essentially constant outer systems. Taking the extra tail in Section 2 preserves their pro-objects,\allowbreak{} so the same middle value exists. At 8348–\allowbreak{}8362 the cited restriction triangle can be rotated,\allowbreak{} with its prescribed rotation sign,\allowbreak{} so any pair of essentially constant terms gives the third; exact restriction/\allowbreak{}direct-image functors preserve the identities (1.1)\allowbreak{}. Thus the categorical inference is unchanged. Earlier geometric and sheaf-theoretic dependencies remain in that chapter\textquotesingle s review.

In more-algebra.tex 27859–\allowbreak{}27874,\allowbreak{} the negative truncations of the \(r\)-term Koszul complexes have cohomology only in the \(r\) original degrees \([-r,-1]\); once their individual cohomology transitions are killed by the subsequent argument,\allowbreak{} \(r\) successive simultaneous choices kill the derived transition. The augmentation triangle then gives the same pro-isomorphism. The zero-generator case has zero negative truncation.

At 29704–\allowbreak{}29724 and 30366–\allowbreak{}30389,\allowbreak{} the receiving proof first replaces an unbounded derived tensor presentation by the pro-isomorphic Koszul presentation with a common finite cohomology range. The criterion is applied to that bounded presentation,\allowbreak{} and pro-isomorphism is then transported back by the actual tensor functor. No boundedness of the original derived tensor objects is inferred. At 30333–\allowbreak{}30352,\allowbreak{} the cone sequence has the finite range coming from the displayed finite complex; its known pro-zero cohomology makes it pro-zero. These are the exact categorical uses checked here. They do not certify the unreviewed Artin–\allowbreak{}Rees inputs or prove the entire receiving chapter.

All mathematical text of the Derived authority chapter has now been read and semantically reviewed through 12590; terminal apparatus through 12604 has been read. The final bibliography/\allowbreak{}input commands contain no additional substantive mathematical result. The next step is to consolidate the chapter\textquotesingle s complete evidence and proposed operations,\allowbreak{} retaining all inherited corrections,\allowbreak{} conceptual errata,\allowbreak{} editorial consequences,\allowbreak{} and receiver followups before admission/\allowbreak{}composition and the required build/\allowbreak{}review/\allowbreak{}publication workflow. Completion of this chapter\textquotesingle s source review does not complete the larger correction queue or the active programme.
